\chapter{Introducción}

La transformación digital en los entornos industriales ha impulsado la necesidad de desarrollar soluciones tecnológicas que optimicen la gestión de la información, mejoren la eficiencia operativa y faciliten la toma de decisiones estratégicas. En este contexto, las organizaciones que dependen de procesos productivos complejos requieren sistemas informáticos capaces de integrarse de manera efectiva, garantizar la disponibilidad de datos en tiempo real y asegurar la trazabilidad de sus operaciones.

El presente documento tiene como propósito describir el desarrollo de una solución tecnológica orientada a fortalecer el Sistema Integral de Producción dentro del entorno industrial de TenarisTamsa. A partir del análisis de la problemática existente, se identifican diversas limitaciones relacionadas con la gestión manual de información, la falta de integración entre sistemas y la necesidad de contar con herramientas que permitan un mejor control y seguimiento de los procesos productivos.

Con base en lo anterior, se plantea el diseño e implementación de una aplicación web que permita digitalizar procesos, centralizar la información y mejorar la interacción de los usuarios con el sistema, utilizando tecnologías modernas de desarrollo tanto en el front-end como en el back-end. Esta propuesta busca no solo atender las necesidades actuales de la organización, sino también sentar las bases para la escalabilidad y mejora continua de sus sistemas informáticos.

Este trabajo refleja la aplicación de conocimientos teóricos y prácticos en un entorno real, contribuyendo al desarrollo de soluciones tecnológicas que impactan directamente en la eficiencia y competitividad de la organización.